% ============================================================
% Presentación 1 — ICS2122 Taller de Investigación Operativa
% Ninguna cifra está escrita a mano: todas vienen de
% ../../03-resultados/, generado por 02-analisis-datos/.
% ============================================================
\documentclass[aspectratio=169,11pt]{beamer}
% ============================================================
% estilo_beamer.tex — preámbulo de la presentación
% Réplica de la paleta de 04-documentos/estilo.tex para que el
% informe y el deck se vean como un mismo trabajo.
% ============================================================

\usepackage[utf8]{inputenc}
\usepackage[T1]{fontenc}
\usepackage[spanish,es-noshorthands]{babel}
\usepackage{lmodern}
\usepackage{booktabs}
\usepackage{array}
\usepackage{colortbl}
\usepackage{tabularx}
\usepackage{graphicx}
\usepackage{xcolor}

% ---------- paleta compartida con el informe ----------
\definecolor{ink}{RGB}{0,62,92}
\definecolor{inklight}{RGB}{92,132,153}
\definecolor{rulegray}{RGB}{205,214,219}
\definecolor{critred}{RGB}{158,38,32}
\definecolor{amber}{RGB}{176,98,16}
\definecolor{slate}{RGB}{104,112,118}
\definecolor{headrow}{RGB}{0,62,92}
\definecolor{softrow}{RGB}{242,246,248}

% ---------- tema ----------
\usetheme{default}
\usecolortheme{default}
\setbeamertemplate{navigation symbols}{}
\setbeamercolor{structure}{fg=ink}
\setbeamercolor{frametitle}{fg=ink}
\setbeamercolor{title}{fg=ink}
\setbeamercolor{itemize item}{fg=inklight}
\setbeamercolor{itemize subitem}{fg=inklight}
\setbeamerfont{frametitle}{series=\bfseries,size=\large}
\setbeamerfont{title}{series=\bfseries}
\setbeamertemplate{itemize item}{\textbf{\textperiodcentered}}

% línea bajo el título de cada lámina
\setbeamertemplate{frametitle}{%
  \vspace{6pt}%
  \usebeamerfont{frametitle}\usebeamercolor[fg]{frametitle}\insertframetitle\par
  \vspace{2pt}{\color{rulegray}\rule{\textwidth}{0.6pt}}\par\vspace{-2pt}%
}

% ---------- numeración de diapositivas (exigida por los lineamientos) ----------
\setbeamertemplate{footline}{%
  \hfill\usebeamerfont{footline}\color{slate}\scriptsize
  \insertframenumber{} / \inserttotalframenumber\hspace{12pt}\vspace{6pt}%
}

% ---------- tablas, igual que en el informe ----------
\newcolumntype{L}[1]{>{\raggedright\arraybackslash}p{#1}}
\newcommand{\thead}[1]{\textcolor{white}{\bfseries #1}}
\renewcommand{\arraystretch}{1.2}
\setlength{\tabcolsep}{4pt}

% ---------- cifra destacada ----------
% #1 = número   #2 = qué es
\newcommand{\cifra}[2]{%
  \begin{minipage}[t]{0.31\textwidth}
    \raggedright
    {\color{ink}\bfseries\fontsize{22}{24}\selectfont #1}\\[2pt]
    {\color{slate}\footnotesize #2}
  \end{minipage}%
}

% variante apilada: ocupa el ancho de su columna
\newcommand{\cifravert}[2]{%
  \par\noindent
  {\color{ink}\bfseries\fontsize{17}{19}\selectfont #1}\ \ {\color{slate}\footnotesize #2}%
  \par\addvspace{6pt}%
}

% ---------- conclusión de la lámina ----------
\newcommand{\remate}[1]{%
  \vspace{4pt}{\color{rulegray}\rule{\textwidth}{0.6pt}}\par\vspace{3pt}
  {\color{ink}\small\textbf{#1}}%
}

% ---------- etiquetas de estado ----------
\newcommand{\etiqueta}[2]{%
  {\setlength{\fboxsep}{2.5pt}\colorbox{#1}{\textcolor{white}{\bfseries\scriptsize\strut\,#2\,}}}%
}


\newcommand{\RES}{../../03-resultados}
\input{\RES/metricas}
\input{\RES/metricas_entrega1}

\title{Planificación de la operación de buses eléctricos}
\subtitle{Red Metropolitana de Movilidad — Entrega 1}
\author{Grupo \textbf{[N]} \quad\textbar\quad ICS2122 Taller de Investigación Operativa}
\date{3 de septiembre de 2026}

\begin{document}

% ------------------------------------------------------------------ 1
\begin{frame}[plain]
  \titlepage
\end{frame}

% ------------------------------------------------------------------ 2
\begin{frame}{El problema}
  \small
  Programar la operación de \textbf{un día completo} de la flota de buses
  eléctricos de la RED: qué bus cubre qué viaje, cómo se encadena con el
  siguiente, y cuándo y dónde recarga.

  \vspace{6pt}
  \begin{columns}[t]
    \begin{column}{0.31\textwidth}
      {\color{ink}\bfseries Qué se decide}
      \begin{itemize}\footnotesize
        \item Asignación de viajes a buses
        \item Encadenamiento y traslados en vacío
        \item Momento, lugar y duración de cada recarga
      \end{itemize}
    \end{column}
    \begin{column}{0.31\textwidth}
      {\color{ink}\bfseries Qué restringe}
      \begin{itemize}\footnotesize
        \item Cada viaje se cubre exactamente una vez
        \item Autonomía: \autonomiaUtil{} km útiles por carga
        \item \numElectroterminales{} electroterminales, \cuposTotales{} cupos
        \item Potencia eléctrica limitada
      \end{itemize}
    \end{column}
    \begin{column}{0.31\textwidth}
      {\color{ink}\bfseries Qué entregamos}
      \begin{itemize}\footnotesize
        \item Un plan diario factible
        \item Evaluado con KPIs de flota, energía, costo y servicio
        \item Con una cota que diga cuán lejos está del óptimo
      \end{itemize}
    \end{column}
  \end{columns}

  \remate{En la literatura se llama \textit{Multi-Depot Electric Vehicle
  Scheduling Problem} con programación de carga acoplada. El enunciado no lo
  nombra, pero esa etiqueta es la llave de toda la bibliografía.}
\end{frame}

% ------------------------------------------------------------------ 3
\begin{frame}{Con qué datos contamos}
  \scriptsize
  \begin{columns}[t]
    \begin{column}{0.56\textwidth}
      \begin{tabular}{@{}L{2.5cm} L{4.4cm}@{}}
        \rowcolor{headrow}
        \thead{Archivo} & \thead{Qué aporta} \\
        \midrule
        \texttt{gtfs/trips} & \registrosViaje{} viajes de bus del día laboral \\
        \rowcolor{softrow}
        \texttt{gtfs/frequencies} & Convierte cada viaje en varias corridas reales \\
        \texttt{gtfs/stop\_times} & Horarios y secuencia de paradas \\
        \rowcolor{softrow}
        \texttt{gtfs/stops} & Coordenadas de \textbf{12.877} paradas \\
        \texttt{shape\_distances} & Distancia de cada recorrido \\
        \rowcolor{softrow}
        \texttt{charger\_capacity} & \numElectroterminales{} electroterminales, \cuposTotales{} cupos \\
        \texttt{electricity\_prices} & Tarifa por tramo horario \\
        \rowcolor{softrow}
        \texttt{parameters} & Batería, autonomía, potencia, costos \\
        \bottomrule
      \end{tabular}
    \end{column}
    \begin{column}{0.42\textwidth}
      {\color{critred}\bfseries\small Lo que NO está}
      \begin{itemize}\scriptsize
        \item \camposAusentes{} de \camposDocumentados{} campos del diccionario
              no existen en ningún archivo: terminales de cada viaje, distancias
              de traslado en vacío, número de cargadores, potencia del
              transformador.
      \end{itemize}
      \vspace{4pt}
      {\color{ink}\bfseries\small Lo que sí pudimos reconstruir}
      \begin{itemize}\scriptsize
        \item Origen y destino de \pctGeoReconstruido{}\,\% de los viajes,
              cruzando \texttt{stop\_times} con \texttt{stops}, y con eso los
              traslados en vacío que faltaban.
      \end{itemize}
    \end{column}
  \end{columns}
  \remate{Un dato ausente no siempre es un dato perdido: buena parte de lo
  que falta se deriva del propio GTFS.}
\end{frame}

% ------------------------------------------------------------------ 4
\begin{frame}{Cómo es un día de operación}
  \begin{columns}[c]
    \begin{column}{0.5\textwidth}
      \includegraphics[width=\linewidth]{\RES/fig_duracion.pdf}\\[2pt]
      \includegraphics[width=\linewidth]{\RES/fig_distancia.pdf}
    \end{column}
    \begin{column}{0.48\textwidth}
      \cifravert{\corridasReales}{corridas que cubrir}
      \cifravert{\duracionMediana{} min}{mediana de viaje}
      \cifravert{\distanciaMediana{} km}{mediana de viaje}
      \cifravert{\kmTotales{} km}{en el día}

      \vspace{1pt}
      {\scriptsize\color{slate}\itshape \registrosViaje{} registros en
      \texttt{trips.txt} no son \registrosViaje{} viajes: se repiten según
      \texttt{frequencies.txt} y expandirlos multiplica el problema por
      \factorFrecuencia{}.}
    \end{column}
  \end{columns}
\end{frame}

% ------------------------------------------------------------------ 5
\begin{frame}{La curva del día fija el tamaño de la flota}
  \begin{columns}[c]
    \begin{column}{0.62\textwidth}
      \includegraphics[width=\linewidth]{\RES/fig_concurrencia.pdf}
    \end{column}
    \begin{column}{0.36\textwidth}
      \scriptsize
      A las \horaPico{} hay \textbf{\picoSimultaneo{} viajes simultáneos}.
      Ningún plan puede usar menos buses que eso: es una cota inferior dura.

      \vspace{6pt}
      Validación externa: la RED operaba \flotaRealRED{} buses en 2018.
      Nuestra cota queda a \desvioValidacion\,\% de esa cifra, lo que sugiere
      que el cálculo es correcto.

      \vspace{6pt}
      {\color{critred}\bfseries El problema:} el archivo de parámetros declara
      \flotaDeclarada{} buses.
    \end{column}
  \end{columns}
  \remate{La forma de la curva, no el promedio, es lo que dimensiona la flota.}
\end{frame}

% ------------------------------------------------------------------ 6
\begin{frame}{Los datos no vienen limpios: \chequeosTotal{} verificaciones}
  \scriptsize
  \input{\RES/tabla_calidad_deck}
  \vspace{4pt}
  \remate{\chequeosOk{} verificaciones en orden, \chequeosAlerta{} alertas y
  \chequeosFalla{} fallas. Las dos fallas cambian el alcance del proyecto,
  no son detalles de limpieza.}
\end{frame}

% ------------------------------------------------------------------ 7
\begin{frame}{Dos hallazgos que obligan a filtrar antes de modelar}
  \small
  \begin{columns}[t]
    \begin{column}{0.48\textwidth}
      {\color{ink}\bfseries 1. El GTFS no trae solo buses}
      \begin{itemize}\footnotesize
        \item \viajesNoBus{} de \viajesDelDia{} viajes del día
              (\pctNoBus\,\%) son líneas de Metro y el Metrotren Rancagua.
        \item Programarlos con la flota de buses eléctricos no tiene sentido:
              se descartan filtrando por \texttt{route\_type}.
        \item Tres paradas del feed están fuera de la Región Metropolitana,
              a más de 90 km del centro: son estaciones del Metrotren.
      \end{itemize}
    \end{column}
    \begin{column}{0.48\textwidth}
      {\color{ink}\bfseries 2. El consumo declarado no calza}
      \begin{itemize}\footnotesize
        \item \texttt{battery\_kwh}\,/\,\texttt{range\_km} implica
              \consumoImplicito{} kWh/km.
        \item El piloto real Metbus--Enel--BYD en Santiago reporta
              \consumoReal{} kWh/km.
        \item La diferencia es de \textbf{\desvioConsumo\,\%}: son \energiaImplicita{} GWh/día
              frente a \energiaReal{}.
      \end{itemize}
    \end{column}
  \end{columns}
  \remate{Sin estos dos filtros, cada cifra del proyecto arrastra el error.
  Trabajamos con el valor del curso y reportamos sensibilidad al otro.}
\end{frame}

% ------------------------------------------------------------------ 8
\begin{frame}{Qué le toca a un bus típico}
  \vspace{4pt}
  \cifra{\viajesPorBus}{viajes por bus al día}\hfill
  \cifra{\kmPorBus~km}{recorridos por bus al día}\hfill
  \cifra{\recargasPorBus}{recargas por bus al día}

  \vspace{16pt}
  \small
  Cada recarga completa toma \tiempoRecarga{} minutos a
  \textbf{180 kW}, y devuelve \autonomiaUtil{} km útiles.
  Un bus consume \energiaPorBus{} kWh en el día: menos de una batería
  completa, pero no lo suficiente como para ignorar la recarga.

  \vspace{6pt}
  La flota entera necesita \energiaImplicita{} GWh diarios. Los
  \cuposTotales{} cupos suman \potenciaInstalada{} MW instalados y podrían
  entregar \capacidadCargaDia{} GWh en 24 h: usamos el
  \usoCapacidadCarga\,\% de esa capacidad.

  \remate{En agregado la energía alcanza. El problema es que nadie recarga
  en agregado: todos quieren hacerlo en la misma ventana nocturna.}
\end{frame}

% ------------------------------------------------------------------ 9
\begin{frame}{La restricción que define el alcance del proyecto}
  \small
  \begin{columns}[t]
    \begin{column}{0.5\textwidth}
      El servicio completo exige \horasBus{} horas-bus.

      \vspace{4pt}
      El pico de viajes simultáneos exige \picoSimultaneo{} buses,
      antes de considerar traslados y recargas.

      \vspace{8pt}
      {\color{critred}\bfseries Los \flotaDeclarada{} buses declarados
      son menos que los viajes simultáneos en la hora punta.}

      \vspace{8pt}
      Con \flotaDeclarada{} buses cada uno tendría que hacer
      \viajesPorBusDecl{} viajes y \kmPorBusDecl{} km diarios: son
      \recargasPorBusDecl{} recargas por bus.
    \end{column}
    \begin{column}{0.46\textwidth}
      {\color{ink}\bfseries Esto no es un problema computacional}

      \vspace{4pt}
      Es infactible por construcción. Solo hay tres lecturas posibles:

      \begin{enumerate}\footnotesize
        \item La instancia es un \textbf{subconjunto} de la red, no la red completa.
        \item Los \flotaDeclarada{} buses son referenciales y la flota es una variable.
        \item Hay un error en los datos entregados.
      \end{enumerate}

      \vspace{6pt}
      \textit{Es la primera pregunta que le haremos al profesor. Mientras tanto
      trabajamos con la lectura (1).}
    \end{column}
  \end{columns}
\end{frame}

% ------------------------------------------------------------------ 10
\begin{frame}{Por qué el problema es difícil}
  \begin{columns}[c]
    \begin{column}{0.44\textwidth}
      \includegraphics[width=\linewidth]{\RES/fig_deadhead.pdf}
      \vspace{2pt}
      {\scriptsize\color{slate} Distancia entre el fin de un viaje y el inicio
      de otro con el que podría encadenarse.}
    \end{column}
    \begin{column}{0.54\textwidth}
      \small
      Encadenar viajes genera \textbf{\arcosFull{} millones} de conexiones
      horariamente posibles: unas \arcosPorCorrida{} por corrida.
      Cada una sería una variable binaria.

      \vspace{6pt}
      Pero la geografía las poda:
      \begin{itemize}\footnotesize
        \item De las conexiones con menos de \holguraAjustada{} min de holgura,
              solo el \fraccionGeoFactible\,\% da tiempo de trasladarse.
        \item Solo el \fraccionBarata\,\% de las conexiones tiene un traslado
              corto (bajo \deadheadBarato{} km).
        \item La mediana del traslado entre dos viajes encadenables es
              \deadheadMediana{} km; el percentil 90 llega a \deadheadPnoventa{} km.
      \end{itemize}
    \end{column}
  \end{columns}
  \remate{Casi todas las conexiones legales son malas. El valor del modelo está
  en encontrar ese \fraccionBarata\,\% que además cierra el balance de energía.}
\end{frame}

% ------------------------------------------------------------------ 11
\begin{frame}{Familia de instancias: de lo verificable a lo real}
  \scriptsize
  \input{\RES/tabla_instancias_deck}
  \vspace{4pt}
  \remate{Todas son el mismo día laboral; cambia cuántas de las \numRutas{}
  rutas se incluyen. La columna de arcos es la que decide qué método es viable
  en cada escala.}
\end{frame}

% ------------------------------------------------------------------ 12
\begin{frame}{Con qué vamos a medir una solución}
  \scriptsize
  \input{\RES/tabla_kpis_deck}
  \vspace{3pt}
  \remate{Los ocho primeros ya tienen valor de referencia sobre los datos
  entregados; los cuatro últimos solo existen una vez que hay un plan.}
\end{frame}

% ------------------------------------------------------------------ 13
\begin{frame}{Qué encontramos en la literatura}
  \small
  \begin{columns}[t]
    \begin{column}{0.32\textwidth}
      {\color{ink}\bfseries El acoplamiento es la novedad}
      \footnotesize\\[3pt]
      La revisión de Perumal et al. (2022) reporta que la literatura resuelve
      el \textit{scheduling} de vehículos y la programación de la recarga
      \textbf{por separado}. El enunciado exige tomarlas juntas.
    \end{column}
    \begin{column}{0.32\textwidth}
      {\color{ink}\bfseries Los exactos no escalan}
      \footnotesize\\[3pt]
      El mejor método exacto publicado (Wu et al., 2025) resuelve a optimalidad
      instancias de \textbf{101 corridas}. Nuestro caso completo tiene
      \corridasReales{}: tres órdenes de magnitud más.
    \end{column}
    \begin{column}{0.32\textwidth}
      {\color{ink}\bfseries La capacidad de carga casi no se modela}
      \footnotesize\\[3pt]
      De Vos et al. (2024) es de los pocos trabajos que combina capacidad
      limitada de estación con recarga parcial, que es exactamente nuestra
      variante.
    \end{column}
  \end{columns}
  \remate{La brecha entre lo que la literatura resuelve exactamente y el tamaño
  de nuestro caso es la que determina la metodología, no una preferencia nuestra.}
\end{frame}

% ------------------------------------------------------------------ 14
\begin{frame}{Metodologías consideradas}
  \scriptsize
  \begin{tabular}{@{}L{2.6cm} L{2.2cm} L{2.9cm} L{4.3cm}@{}}
    \rowcolor{headrow}
    \thead{Familia} & \thead{Garantía} & \thead{Escala reportada} & \thead{Por qué no basta sola} \\
    \midrule
    Flujo en \textit{time-space} & Óptimo del modelo & Instancias reales & Al incorporar el estado de carga deja de ser flujo puro \\
    \rowcolor{softrow}
    \textit{Branch-and-price} & Óptimo con cota & 101 corridas & Tres órdenes de magnitud bajo nuestro caso \\
    Metaheurísticas (ALNS) & Sin garantía & Instancias grandes & Sin cota, no se puede defender la calidad \\
    \rowcolor{softrow}
    Descomposición de Benders & Óptimo o brecha $<$1\,\% & Datos de Chicago & Supone decisiones de inversión que aquí ya están dadas \\
    Modelos estocásticos & Probabilística & No reportada & Exigen distribuciones que los datos no traen \\
    \rowcolor{softrow}
    Simulación & No optimiza & --- & Mide fragilidad, pero no la corrige \\
    \bottomrule
  \end{tabular}
  \remate{Ninguna familia cubre sola el caso: exacta y escala son excluyentes
  a este tamaño.}
\end{frame}

% ------------------------------------------------------------------ 15
\begin{frame}{Ruta escogida: híbrida en dos etapas}
  \small
  \begin{columns}[t]
    \begin{column}{0.3\textwidth}
      {\color{ink}\bfseries Etapa 1 — Exacta}
      \footnotesize\\[3pt]
      Modelo entero sobre \textit{time-space network} con la capa energética,
      resuelto sobre \textbf{I-001} e \textbf{I-010}.
      \\[3pt]
      \textit{Para qué:} verificar que el modelo representa el problema y
      producir la cota de referencia.
    \end{column}
    \begin{column}{0.3\textwidth}
      {\color{ink}\bfseries Etapa 2 — Heurística}
      \footnotesize\\[3pt]
      ALNS extendido con la restricción de capacidad de carga, aplicado a
      \textbf{I-050} y hacia arriba.
      \\[3pt]
      \textit{Para qué:} llegar a la escala real. Su calidad se mide contra el
      óptimo de la Etapa 1.
    \end{column}
    \begin{column}{0.34\textwidth}
      {\color{ink}\bfseries Complemento — Robustez}
      \footnotesize\\[3pt]
      Simulación sobre el plan final y análisis de sensibilidad en consumo
      (\consumoReal{} vs. \consumoImplicito{} kWh/km) y en tiempos de traslado.
      \\[3pt]
      \textit{Por qué así:} un modelo estocástico completo exige distribuciones
      que los datos no contienen.
    \end{column}
  \end{columns}
  \remate{Lo que perdemos: la solución del caso grande no tendrá garantía de
  optimalidad. Lo decimos de entrada, y por eso la Etapa 1 existe.}
\end{frame}

% ------------------------------------------------------------------ 16
\begin{frame}{Pasos siguientes}
  \small
  \begin{columns}[t]
    \begin{column}{0.5\textwidth}
      {\color{ink}\bfseries Próximas cuatro semanas}
      \begin{itemize}\footnotesize
        \item Validar los supuestos abiertos con el equipo docente
              (flota, consumo, capacidad de carga, deadhead).
        \item Construir el generador de instancias y congelar \textbf{I-010}.
        \item Formular y resolver el modelo exacto sobre \textbf{I-001}.
        \item Primer plan factible sobre \textbf{I-010} con KPIs medidos.
      \end{itemize}
    \end{column}
    \begin{column}{0.46\textwidth}
      {\color{ink}\bfseries Lo que ya está resuelto}
      \begin{itemize}\footnotesize
        \item Datos verificados, filtrados y reproducibles de punta a punta.
        \item Geografía reconstruida: los traslados en vacío ya son calculables.
        \item Familia de instancias definida y dimensionada.
        \item KPIs y metodología escogida y argumentada.
      \end{itemize}
    \end{column}
  \end{columns}
  \remate{La pregunta del principio era qué bus cubre qué viaje y cuándo
  recarga. Hoy sabemos cuán grande es esa pregunta, qué datos la responden y
  cuáles todavía faltan.}
\end{frame}

% ------------------------------------------------------------------ anexo
\appendix
\begin{frame}{Anexo — Tabla de instancias completa}
  \tiny
  \input{\RES/tabla_instancias}
\end{frame}

\begin{frame}{Anexo — Verificación de datos completa}
  \tiny
  \input{\RES/tabla_calidad}
\end{frame}

\begin{frame}{Anexo — Distancia a electroterminal}
  \begin{columns}[c]
    \begin{column}{0.6\textwidth}
      \includegraphics[width=\linewidth]{\RES/fig_pullout.pdf}
    \end{column}
    \begin{column}{0.38\textwidth}
      \scriptsize
      Distancia desde el inicio de cada viaje al electroterminal más cercano.
      Media \pullOutMedio{} km, percentil 90 \pullOutPnoventa{} km.
      \\[4pt]
      Supuestos: velocidad en vacío \velocidadDeadhead{} km/h,
      factor de rodeo \factorRodeo{} sobre la distancia recta.
      Velocidad comercial observada: \velocidadComercial{} km/h.
    \end{column}
  \end{columns}
\end{frame}

\end{document}
